\documentclass{article}

\usepackage[preprint]{neurips_2025}
\usepackage[utf8]{inputenc}
\usepackage[T1]{fontenc}
\usepackage{hyperref}
\usepackage{url}
\usepackage{booktabs}
\usepackage{amsmath}
\usepackage{amsfonts}
\usepackage{microtype}

\setlength{\tabcolsep}{4pt}
\renewcommand{\arraystretch}{0.92}
\setlength{\textfloatsep}{7pt}
\setlength{\floatsep}{6pt}
\setlength{\intextsep}{6pt}

\title{Layer-Adaptive Pyramid SinkKV for Training-Free KV Cache Compression}

\author{
Lixiang Wang\\
\texttt{cort\_lee@sjtu.edu.cn}
\And
Rui Cai\\
\texttt{cai\_rui@sjtu.edu.cn}
\And
Haotian Wu\\
\texttt{wuhaotian@sjtu.edu.cn}
}

\begin{document}

\maketitle

\begin{abstract}
Long-context decoding with Transformer language models is increasingly limited by
the memory footprint and attention cost of the key-value (KV) cache. This report
studies training-free KV cache compression for efficient inference under the
course constraint that model parameters remain unchanged. Building on attention
sinks, attention-selected middle tokens, and recent-window retention, we propose
\textbf{Layer-Adaptive Pyramid SinkKV}, a layer-wise cache pruning policy that
allocates larger budgets to lower layers and smaller budgets to higher layers.
The method keeps the same average nominal budget as a uniform Sink-SnapKV
baseline while changing where the retained cache is placed across depth. On
Pythia-70M, Pyramid SinkKV improves Wikitext-2 perplexity from 35.68 to 35.05
relative to the uniform baseline with a slightly smaller measured average KV
cache, but it is slightly worse on the tested PG-19 sample, increasing
perplexity from 31.23 to 31.46. A reverse-pyramid ablation is worse on both
datasets, suggesting that layer direction matters. These results do not show a
universal win, but they provide evidence that layer-wise budget allocation is a
relevant design factor for training-free KV cache compression.
\end{abstract}

\section{Introduction}
\label{sec:intro}

Autoregressive language model inference stores a KV cache for all previously
processed tokens. The cache avoids recomputing past keys and values, but its
sequence length grows with the context and directly affects GPU memory usage and
the attention work performed during decoding. For long prompts and streaming
generation, a dense cache can become a practical bottleneck even when the model
itself is small enough to fit in memory. The assignment therefore asks for
training-free inference optimization on Pythia-70M, with quality and efficiency
measured on datasets such as Wikitext-2 and PG-19.

The project implementation is available at
\url{https://github.com/presucc/kv-cache-compression}. It reproduces several
cache retention policies and implements the proposed Pyramid SinkKV variant.
The individual results show the central trade-off: aggressive pruning
can reduce retained tokens or improve time per output token, but it can also
discard context that is important for perplexity. This motivates evaluating not
only how many tokens are compressed, but also how the retained budget is
structured. A cache policy should specify which tokens are retained, why the
retained budget is allocated that way, and where the quality loss comes from.

We focus on the layer allocation problem. Existing simple baselines often apply
the same KV budget to every Transformer layer. However, different layers can
play different roles: lower layers are closer to lexical and local-context
features, while higher layers tend to aggregate information into more compact
representations. This motivates a pyramid allocation that keeps more recent and
selected historical tokens in lower layers and prunes more aggressively in upper
layers. The resulting method, Layer-Adaptive Pyramid SinkKV, is still
training-free and does not modify model weights.

Our contributions are threefold. First, we formulate a compact layer-adaptive
extension of Sink-SnapKV that combines sink tokens, attention-selected middle
tokens, and a recent window under a pyramid budget. Second, we evaluate dense
attention, StreamingLLM, uniform Sink-SnapKV, Pyramid SinkKV, and a
reverse-pyramid ablation under the same Pythia-70M pipeline. Third, we report a
mixed empirical finding: pyramid allocation improves Wikitext-2 quality and
beats the reverse ablation, while PG-19 exposes that a fixed pyramid is not
always better than a uniform middle-context budget.

\section{Related Work}
\label{sec:related-work}

\paragraph{StreamingLLM.}
Long-context generation is bottlenecked by the growth of the KV cache: every
decoded token appends new key and value states, so memory and attention cost
increase with the effective context length. A direct way to control this cost is
a sliding-window cache that keeps only the most recent tokens. However,
StreamingLLM observes that pure window attention can fail once early
``attention sink'' tokens are evicted, because many language models assign
disproportionate attention mass to a small set of initial positions even when
their semantic content is not directly relevant \citep{xiao2023streamingllm}.
StreamingLLM therefore keeps a small number of sink tokens together with a recent
window. This design does not explicitly preserve semantically important
middle-context tokens.

\paragraph{SnapKV.}
SnapKV addresses this limitation by using attention information from the prompt
to predict which historical KV entries will be useful during subsequent
generation \citep{li2024snapkv}. It uses an observation window near the end of
the prompt, aggregates attention weights over earlier positions, smooths the
importance scores, and retains the top-scoring KV positions for each head. This
design is content-adaptive and training-free, but obtaining and processing
attention scores can reduce realized speedup, especially on small models or
short prompts where KV attention is not the dominant cost.

\paragraph{PyramidKV.}
PyramidKV argues that KV compression should not use a uniform budget across
layers \citep{cai2024pyramidkv}. Its key observation is pyramidal information
funneling: lower layers distribute attention more broadly across the context,
while higher layers aggregate information toward fewer critical positions. A
layer-wise compression policy can therefore allocate larger budgets to lower
layers and smaller budgets to higher layers. Our method adapts this idea to a
small course-project setting by combining it with sink tokens, a recent window,
and attention-selected middle-context tokens.

These three lines are complementary: StreamingLLM contributes robust sink
tokens, SnapKV contributes content-aware selection, and PyramidKV contributes
layer-wise budgeting. Pyramid SinkKV combines these ideas while keeping the
inference-only nature of the baselines.

\section{Methodology}
\label{sec:methodology}

We introduce \textbf{Layer-Adaptive Pyramid SinkKV}, a training-free KV cache
compression method for decoder-only language models. The method preserves the
same three cache components as the uniform Sink-SnapKV baseline: initial sink
tokens, a recent local window, and attention-selected middle-context tokens. Its
main change is that the sizes of the recent and selected components depend on
the Transformer layer.

For layer $\ell$, the retained cache positions are
\[
K_\ell = S \cup R_\ell \cup I_\ell ,
\]
where $S$ is the shared set of sink tokens, $R_\ell$ is the layer-specific recent
window, and $I_\ell$ contains important middle tokens selected from attention
scores. The model weights are unchanged, so the method only modifies
inference-time cache pruning. In the implementation, cache pruning is performed
after decoding steps using a separate index set for each layer. When attention
scores are unavailable, the selected-token component degenerates to a recent
fallback; the reported attention-based experiments therefore use eager attention
so that token importance can be computed explicitly.

For the 6-layer Pythia-70M model, the budget profile is fixed before inference:
layers 0--1 keep $(4,384,48)$ sink, window, and selected tokens; layers 2--3 keep
$(4,256,32)$; and layers 4--5 keep $(4,128,16)$. The average nominal budget is
292 tokens per layer, exactly matching the uniform Sink-SnapKV baseline
$(4+256+32=292)$. Thus the comparison tests the effect of budget redistribution,
not a larger total cache. We also evaluate \texttt{reverse\_pyramid}, which
swaps the low-layer and high-layer budgets, to test whether the allocation
direction matters.

Within each layer, sink tokens stabilize attention, recent tokens preserve local
coherence, and selected middle tokens retain non-local evidence. After these
sets are merged, the KV tensors of that layer are indexed by the corresponding
keep positions. We use a fixed schedule rather than a dynamic controller so that
the experiment isolates the effect of layer-wise allocation.

\section{Experiments}
\label{sec:experiments}

\paragraph{Setup.}
We evaluate \texttt{EleutherAI/pythia-70m} on Wikitext-2 validation text and one
long PG-19 test sample. PPL is measured with 1024 evaluation tokens. The PPL
runs use CUDA with \texttt{float32}; generation latency uses CUDA with
\texttt{float16}, a 512-token prompt, and 64 generated tokens. We report PPL,
maximum retained KV tokens, average retained KV tokens, time to first token
(TTFT), time per output token (TPOT), throughput, and peak CUDA memory. The
average retained KV count is averaged across decoding steps and, for
layer-adaptive methods, across layers. This is why Pyramid SinkKV can have a
higher maximum retained token count than uniform Sink-SnapKV while still using a
slightly smaller measured average cache.

All compared methods use the same tokenizer, checkpoint, input text, and
evaluation loop. The preliminary reproduction and final experiments are not
treated as one shared benchmark because their sequence lengths and latency
scripts differ; the reproduction is included only as motivation for moving
beyond direct uniform compression.

\paragraph{Preliminary baseline reproduction.}
Before implementing the final layer-adaptive method, the project reproduced a
SnapKV-style baseline on Pythia-70M without changing model weights. These runs
used 2048-token PPL sequences, where the first half built and compressed the KV
cache and the second half was evaluated autoregressively. Generation used a
1024-token prompt and 128 generated tokens on an NVIDIA GeForce RTX 4080 Laptop
GPU. The reproduced SnapKV setting used \texttt{max\_capacity=768},
\texttt{observation\_window=64}, and \texttt{kernel\_size=5}.

\begin{table}[t]
  \centering
  \scriptsize
  \caption{Preliminary SnapKV reproduction motivating the final design. KV memory
  is reported in MB, TTFT in seconds, and TPOT in milliseconds.}
  \label{tab:baseline-repro}
  \setlength{\tabcolsep}{3.3pt}
  \begin{tabular}{llrrrrr}
    \toprule
    Setting & Method & PPL & KV (MB) & TTFT (s) & TPOT (ms) & Tok/s \\
    \midrule
    Wiki PPL & Dense & 39.67 & 12.00 & -- & -- & -- \\
    Wiki PPL & SnapKV & 59.00 & 9.75 & -- & -- & -- \\
    PG PPL & Dense & 28.21 & 12.00 & -- & -- & -- \\
    PG PPL & SnapKV & 44.97 & 9.75 & -- & -- & -- \\
    Generation & Dense & -- & -- & 22.76 & 4.94 & 202.53 \\
    Generation & SnapKV & -- & -- & 21.95 & 4.71 & 212.18 \\
    \bottomrule
  \end{tabular}
\end{table}

Table~\ref{tab:baseline-repro} shows the original quality--efficiency trade-off.
SnapKV reduced KV memory from 12.00 MB to 9.75 MB and improved TPOT from 4.94 ms
to 4.71 ms, while PPL increased on both datasets. This diagnostic
result motivated a method that keeps the same three-part cache structure while
testing whether layer-wise budget redistribution can improve the
quality--compression frontier.

\paragraph{Main results and analysis.}

\begin{table}[t]
  \centering
  \small
  \caption{Main PPL results. Max and Avg denote retained KV tokens.}
  \label{tab:main}
  \begin{tabular}{llrrr}
    \toprule
    Data & Method & PPL & Max & Avg \\
    \midrule
    PG-19 & dense & 31.10 & 1023 & 512.00 \\
    PG-19 & streamingllm & 31.43 & 260 & 227.09 \\
    PG-19 & sink\_snapkv & 31.23 & 292 & 250.47 \\
    PG-19 & pyramid\_sinkkv & 31.46 & 436 & 243.71 \\
    \midrule
    WikiText & dense & 30.15 & 1023 & 512.00 \\
    WikiText & streamingllm & 36.19 & 260 & 227.09 \\
    WikiText & sink\_snapkv & 35.68 & 292 & 250.47 \\
    WikiText & pyramid\_sinkkv & 35.05 & 436 & 243.71 \\
    \bottomrule
  \end{tabular}
\end{table}

Table~\ref{tab:main} reports the main comparison. On Wikitext-2, Pyramid SinkKV
improves over uniform Sink-SnapKV, reducing PPL from 35.68 to 35.05 while
retaining a slightly smaller average number of KV tokens. On PG-19, however, it
is slightly worse than uniform Sink-SnapKV, increasing PPL from 31.23 to 31.46.
These results indicate that redistributing the same average budget across layers
can affect the quality--compression trade-off, although the optimal allocation
depends on the dataset and context distribution. PG-19 contains long-form book
text, where preserving a stable uniform middle-context budget may be more
important, whereas Wikitext-2 appears to benefit from allocating additional
capacity to lower layers.

The Wikitext-2 result provides the main positive evidence: Pyramid SinkKV keeps
the same nominal budget as uniform Sink-SnapKV but improves PPL. The PG-19 result
suggests that long-form book text may require a more stable middle-context
budget across layers. Thus a single fixed pyramid schedule is not optimal for
every domain.

\begin{table}[t]
  \centering
  \small
  \caption{Layer-allocation ablation.}
  \label{tab:ablation}
  \begin{tabular}{lrrrr}
    \toprule
    Method & PG PPL & Wiki PPL & Max & Avg \\
    \midrule
    sink\_snapkv & 31.23 & 35.68 & 292 & 250.47 \\
    pyramid\_sinkkv & 31.46 & 35.05 & 436 & 243.71 \\
    reverse\_pyramid & 31.51 & 35.14 & 436 & 243.71 \\
    \bottomrule
  \end{tabular}
\end{table}

Table~\ref{tab:ablation} tests whether the pyramid direction matters. The
reverse-pyramid ablation has worse PPL than Pyramid SinkKV on both datasets,
suggesting that giving larger budgets to lower layers is not arbitrary. The
margin is modest, which is consistent with the limited depth of the 6-layer
Pythia-70M model, but the ordering remains consistent.

\begin{table}[t]
  \centering
  \small
  \caption{PG-19 generation latency with a 512-token prompt and 64 generated
  tokens. TTFT and TPOT are reported in seconds; memory is MB.}
  \label{tab:latency}
  \begin{tabular}{lrrrr}
    \toprule
    Method & TTFT & TPOT & Tok/s & Mem. \\
    \midrule
    dense & 3.85 & 0.0080 & 14.69 & 165.62 \\
    streamingllm & 4.01 & 0.0088 & 14.02 & 154.56 \\
    sink\_snapkv & 4.75 & 0.0092 & 12.01 & 155.71 \\
    pyramid\_sinkkv & 3.81 & 0.0082 & 14.79 & 224.83 \\
    \bottomrule
  \end{tabular}
\end{table}

Table~\ref{tab:latency} reports the latency side. Pyramid SinkKV has comparable
speed and lower measured average KV tokens than uniform Sink-SnapKV, but higher
peak memory. This is likely an implementation artifact of per-layer pruning and
attention bookkeeping rather than an inherent property of the cache policy. The
latency script uses a clear one-token-at-a-time Python loop, so these numbers
should be read as reproducible policy comparisons, not optimized serving
benchmarks.

Overall, the empirical evidence supports a narrower claim than universal
dominance over uniform cache policies: layer-wise cache budgeting is a relevant
design factor. The method improves Wikitext-2 quality at the same nominal budget
and consistently outperforms the reverse allocation, while PG-19 shows that the
best budget profile may depend on long-document structure. A natural next step
is to tune the pyramid profile, reduce attention bookkeeping, and evaluate
whether the same pattern becomes stronger on deeper models.

\begin{thebibliography}{9}

\bibitem[Xiao et~al.(2023)Xiao, Tian, Chen, Han, and Lewis]{xiao2023streamingllm}
Guangxuan Xiao, Yuandong Tian, Beidi Chen, Song Han, and Mike Lewis.
\newblock Efficient streaming language models with attention sinks.
\newblock \emph{arXiv preprint arXiv:2309.17453}, 2023.
\newblock URL: \url{https://arxiv.org/abs/2309.17453}.

\bibitem[Li et~al.(2024)Li, Huang, Yang, Venkitesh, Locatelli, Ye, Cai, Lewis, and Chen]{li2024snapkv}
Yuhong Li, Yingbing Huang, Bowen Yang, Bharat Venkitesh, Acyr Locatelli,
Hanchen Ye, Tianle Cai, Patrick Lewis, and Deming Chen.
\newblock SnapKV: LLM knows what you are looking for before generation.
\newblock \emph{arXiv preprint arXiv:2404.14469}, 2024.
\newblock URL: \url{https://arxiv.org/abs/2404.14469}.

\bibitem[Cai et~al.(2024)Cai, Zhang, Gao, Liu, Li, Liu, Lu, Xiong, Dong, Hu, and Xiao]{cai2024pyramidkv}
Zefan Cai, Yichi Zhang, Bofei Gao, Yuliang Liu, Yucheng Li, Tianyu Liu, Keming
Lu, Wayne Xiong, Yue Dong, Junjie Hu, and Wen Xiao.
\newblock PyramidKV: Dynamic KV cache compression based on pyramidal
information funneling.
\newblock \emph{arXiv preprint arXiv:2406.02069}, 2024.
\newblock URL: \url{https://arxiv.org/abs/2406.02069}.

\end{thebibliography}

\appendix

\section{Work Division and Code Repositories}
\label{app:work-division}

\paragraph{Group work division.}
Lixiang Wang was responsible for the initial idea proposal, the overall paper
structure, and the Introduction. Rui Cai was responsible for the concrete
experiments and the experimental section. Haotian Wu was responsible for idea
refinement, related-literature search, and the Related Work section.

\paragraph{Individual homework repositories.}
The individual code repositories used for the personal part of the assignment
are listed below:
\begin{itemize}
  \item \url{https://github.com/wuhaotian0508/nlp-homework}
  \item \url{https://github.com/presucc/kv-cache-compression}
  \item \url{https://github.com/sjtu-Cort-Lee/NLP-HW-Personal-Part}
\end{itemize}

\end{document}
